\section{Lean definitions, checked results, and the completion path}
\label{sec:lean}

The archive contains a standalone Lean 4.32.0 development using \code{Std}
only. Ten modules were compiled, including the concrete definitions of
all main new programs and their unproved targets. Eleven selected theorem
declarations were audited with \texttt{\#print axioms}. Their dependencies
are contained in \code{propext} and \code{Quot.sound}; none uses
\code{sorryAx} or \code{Classical.choice}. These are facts about the
supplied Lean development, not a claim that all the report's application
theorems have already been formalized.

\subsection{The actual object-language backend}

The namespace \code{LeanProofs.ConciseIPRA} is the repository's namespace.
The reused standalone \code{ObjectKernel.lean} contains the actual
\code{Syntax}, \code{Deduction}, and \code{PR} declarations with their
imports arranged for \code{Std}. The bodies were compared with the
pinned current upstream sources after removing import lines and
boundary whitespace; the source audit records the hashes and equality
checks. The original MIT-0 license is included.

The central types and constructor interfaces are
\begin{lstlisting}
PR : Nat -> Type
PTm : Nat -> Type
Fm PTm n
PRAAx : Theory PTm
Prf PRAAx Gamma phi
\end{lstlisting}
Here \(n\) is the number of free variables. \code{PR k} has constructors
\code{zero}, \code{succ}, \code{proj}, \code{comp}, and \code{prec}.
The last has input symbols of arities \(k\) and \(k+2\), and returns
a symbol of arity \(k+1\). The step's arguments are the recursion index,
the accumulator, and the \(k\) parameters.

The formulas in \code{PRAAx} include the defining equations, successor
separation and injectivity, and quantifier-free induction. The supplied
\code{PRA.qfInd} rule expects a proof that the formula is
\code{Fm.QF}, a base proof, and a step proof in the correctly shifted
context. Equality and substitution are handled by
\code{Prf.eqCongr}, \code{Prf.eqTrans}, \code{Prf.subst}, and
\code{Prf.trans}. The code therefore uses the real proof relation
instead of introducing a new predicate named ``provable'' and assuming
its soundness~\cite{repo-pr}.

\subsection{Status at declaration level}

\begin{longtable}{>{\raggedright\arraybackslash}p{0.27\textwidth}>{\raggedright\arraybackslash}p{0.36\textwidth}>{\raggedright\arraybackslash}p{0.28\textwidth}}
\toprule
Module & Supplied result & Status and logical level\\
\midrule\endhead
\code{ObjectKernel}
& Arithmetic syntax and intuitionistic deduction
& Reused backend; compiled.\\
\code{BackwardDemand}
& Horizon, replay, and finite-query lemmas
& Reused unchanged; checked host theorems.\\
\code{TermCompiler}
& \code{symbol_correct}, \code{transport_equation}
& New checked \code{Prf PRAAx} constructions.\\
\code{DemandCompiler}
& \code{compiled : PR 2}; \code{denote_compiled}; recurrence laws
& New checked syntax and host semantic proofs.\\
\code{ObjectDemand}
& Generic backward zero and successor equations
& New checked \code{Prf PRAAx} constructions.\\
\code{ObjectTargets}
& Internal monus, horizon, replay, and full recurrence statements
& Compiled definitions of unproved propositions.\\
\code{LexFuel}
& PR fuel syntax, including a variable input parameter; checked denotation
& Syntax/semantics checked. Chain and termination theorems are targets.\\
\code{RationalRoots}
& Complete-root candidate program with zero-polynomial failure flag
& Executable definitions; correctness is an unproved target.\\
\code{EliminationTargets}
& Shear, Sylvester determinant, intersection and auxiliary-envelope programs
& Executable definitions; fourteen unproved target propositions.\\
\code{AxiomAudit}
& Dependency reports for eleven selected proved declarations
& Checked audit, not an additional mathematical theorem.\\
\bottomrule
\end{longtable}

All unproved targets are \code{def ... : Prop := ...}, not \code{axiom}
declarations or theorems closed by \code{sorry}. Type checking a target
verifies the well-formedness of its statement and definitions; it does
not prove that statement. The new program definitions can consequently
be inspected and tested without confusing an unfinished theorem with
a trusted assumption.

\subsection{Completing the internal demand theorem}

The first new compiler has an entirely internal correctness conclusion:
\begin{lstlisting}
symbol_correct (t : PTm k) (sigma : Fin k -> PTm n) :
  Prf PRAAx Gamma
    (Fm.eq (PTm.app (symbol t) sigma) (Syn.subst sigma t))
\end{lstlisting}
This theorem is already proved. The two-pass compiler is also a
concrete syntax constructor, but its complete internal correctness
still requires the following named targets.

\begin{enumerate}
\item \code{MonusSelfTarget}. Derive the predecessor and truncated
subtraction equations from \code{PRA.precZero} and
\code{PRA.precSucc}. Use equation induction to prove
\(N\dotminus N=0\) and the elementary successor cancellation laws.
\item \code{BackwardCoherenceTarget}. Use the already checked
\code{object_backward_zero} and \code{object_backward_succ} to prove
\eqref{eq:back-coherent} by QF induction on the reverse index. The
base and step have no order side condition.
\item \code{HorizonFinalTarget}. Substitute the monus-self equation
and the backward zero equation. Define a PR Boolean comparison symbol
and prove the guarded horizon-coherence equation for \(t\le N\).
\item \code{ForwardBaseTarget} and \code{ForwardStepTarget}. Unfold
the exact \code{PR.prec} symbols. Apply the term compiler to the
composite step body, just as in the completed backward proof.
\item Compare the two forward runs with the QF invariant
\(t\le N\to V(t;N+1,K)=V(t;N,D(N,K))\). This is the internal
replacement for the host theorem \code{advance_congr}.
\item \code{CompiledZeroTarget} and \code{CompiledSuccessorTarget}.
Unfold the final composition and apply the preceding equations.
The full human proof appears in Theorem~\ref{thm:demand}.
\end{enumerate}

Every step has a defined numerical term and an exact object formula.
There is no remaining appeal to an arbitrary host induction theorem or
an uninterpreted ``compiler correctness'' axiom. What remains is the
construction of the corresponding finite proof terms.

\subsection{Fixed-dimensional fuel: the uniformity distinction in types}

In \code{LeanProofs.ConciseIPRA.LexFuel},
\code{Rank d := Fin d -> Nat}; \code{lexDropB} implements the
lexicographic comparison. \code{run}, \code{coordMax},
\code{prefixMax}, and \code{trajectoryEnvelope} implement the finite
simulation and the envelope used in the human proof. The host
\code{fuel} definition is exactly \(F_d\).

The checked syntax constructors have types
\begin{lstlisting}
fuelSymbol      : Nat -> PR 1 -> PR 1
fuelParamSymbol : Nat -> PR 2 -> PR 2
boundSymbol     : Nat -> PR 2 -> PR 1
\end{lstlisting}
In the parameterized version, the input symbol \(h\) denotes
\(h(x,t)\), and the output denotes \(F_d(x,s)\). The checked theorem
\code{denote_boundSymbol} proves
\[
 \operatorname{evalB}(\code{boundSymbol}\ d\ h,x)
       =F_d(x,0).
\]
The input \(x\) remains a variable of one PR symbol for every fixed
\(d\); it has not been specialized separately to each numeral.
On the other hand, these constructors do not assert that evaluation of
their generated syntax is a PR function jointly in arbitrary \(d,x,s\).
The non-PR diagonal proved in Section~\ref{sec:descent} explains why.

The exact remaining targets are \code{FuelGrowthTarget},
\code{ControlledChainBoundTarget}, \code{BoundedTerminationTarget},
and \code{WitnessTarget}. Their proof order is:

\begin{enumerate}
\item Prove iterator monotonicity in its starting point and iteration
count. Then prove the fuel growth and monotonicity properties.
\item Implement bounded block boundaries and slices of a finite
controlled chain. The first coordinate drops at each boundary;
within a block apply the preceding-dimensional suffix theorem.
Accumulate the endpoint inequality by induction on the finite number
of blocks.
\item Prove initialization and preservation of the given Boolean
invariant along \code{run}, and persistence of inactivity for the
absorbing transition. An active state at time \(B\) yields the
forbidden \(B+1\)-state chain.
\item Apply the supplied exit implication to the resulting inactive
state. For the internal version, encode tuple and list data by
natural numbers and Booleanize every finite condition. Supply actual
object proofs of the particular transition's local assumptions.
\end{enumerate}

The host \code{Machine} structure deliberately contains data fields,
not unproved correctness axioms. Its arbitrary host functions must
still be replaced by fixed PR syntax before the internal theorem is
instantiated. In particular, well-foundedness of an abstract Lean
relation is not being used as an arithmetic termination certificate.

\subsection{Exact finite polynomial data and executable definitions}

The algebraic specification layer uses
\begin{lstlisting}
QPoly     := List Rat
PlanePoly := List (List Rat)
QPoint    := Rat x Rat
\end{lstlisting}
Coefficients increase with degree; row \(i\), column \(j\) is the
coefficient of \(X^iY^j\). Trailing zeros are permitted. Default reads
return zero. \code{eval} and \code{evalPlane} are Horner evaluation.
The arithmetic layer will encode these finite lists and canonical
rationals by natural numbers; the report's PR constructions provide
the required fixed programs.

\code{RationalRoots.lean} includes the actual denominator product,
integer coefficients, first/last nonzero coefficient selection,
bounded numerator and denominator boxes, evaluation filter, and
duplicate removal. Zero is included as a candidate and tested, which
handles initial powers of the variable. The zero polynomial returns
\code{none}; a nonzero polynomial with no rational roots returns
\code{some []}. Its unproved targets are
\code{RationalRootsTarget}, \code{DegreeBoundTarget}, and
\code{HeightBoundTarget}.

\code{EliminationTargets.lean} supplies actual programs for
\code{shear}, \code{goodShear}, \code{canonicalShear},
\code{sylvester}, \code{determinant}, \code{elimination},
\code{good}, \code{intersect}, \code{auxiliaryBox}, and
\code{envelopeList}. The determinant enumerates the \(N^N\) vectors
of length \(N\) over \(\{0,\ldots,N-1\}\), retains permutations,
computes their signs, and sums their polynomial products. This
inefficient but transparent definition makes its finite bound explicit.
The auxiliary box similarly enumerates finite coefficient vectors.
No polynomial operation in the target module is an opaque oracle.

The mathematical specifications are literal propositions, for example
\begin{lstlisting}
Good F G := good F G = true
common F G p := eval2 F p.1 p.2 = 0 /\ eval2 G p.1 p.2 = 0
\end{lstlisting}
and
\begin{lstlisting}
IntersectCompleteTarget :=
  forall F G, Good F G -> exists points,
    intersect F G = some points /\ points.Nodup /\
    (forall p, p in points <-> common F G p)
\end{lstlisting}
The displayed snippets use ASCII mathematical notation for readability;
the accompanying Lean source contains the checked syntax. The target is
not supplied as a theorem.

\subsection{Dependency order for the plane theorems}

A complete formalization should follow the mathematical dependencies,
which allow the complete-list theorem to be finished before the sharp
count theorem.

\begin{enumerate}
\item \textbf{Rational arithmetic and polynomial identities.}
Prove denominator clearing preserves evaluation, gcd cancellation,
the finite rational-root bound, and correctness of the root scan.
Develop coefficientwise addition, multiplication, evaluation, degree
support, and synthetic division for the finite arrays.
\item \textbf{Shear correctness.}
Prove the binomial coefficient identity for \code{shearCoeff}, then
\code{ShearEvalTarget}. A nonzero polynomial of degree at most
\(d+e\) cannot vanish at all \(d+e+1\) candidate directions; this
proves \code{CanonicalShearTarget} and constant leading coefficients.
\item \textbf{Elimination and complete lists.}
Prove determinant specialization and the adjugate identity. Apply
the evaluation vector \((1,v,\ldots,v^{N-1})\) to establish
\code{CommonGivesRootTarget}. Every specialized first equation has
the same nonzero leading coefficient, so its root routine succeeds.
These facts and the two finite membership scans prove
\code{IntersectCompleteTarget}, followed by
\code{IntersectHeightTarget}.
\item \textbf{Sharp cardinality.}
Prove \code{ResultantDegreeTarget} using the permutation weight
identity from Lemma~\ref{lem:elim-degree}. Prove the Vandermonde
determinant and column factor extraction to obtain
\code{FibreDivisibilityTarget}. Combine distinct linear-factor
powers and apply the degree bound, giving
\code{IntersectCardinalityTarget}.
\item \textbf{Auxiliary exhaustion.}
Prove \code{MonomialCountTarget}, \code{AuxiliaryBoxCountTarget},
and \code{AuxiliaryBoxCoverageTarget}. The last permits a padded
coefficient array different from the input array but coefficientwise
equal to it. Establish invariance of the polynomial programs under
that equality. Finite union membership and the supplied coverage
premise then give \code{EnvelopeCompleteTarget}; the length and
maximum lemmas give \code{EnvelopeCardinalityTarget} and
\code{EnvelopeHeightTarget}.
\end{enumerate}

For the embedded arithmetic layer, every polynomial identity becomes
a bounded conjunction over its finite coefficient support. The support
bound is produced by the same array operations. Every divisibility
argument carries its explicit quotient array. Rational cancellation
carries gcd or B\'ezout data when needed. In the enumeration proofs,
fix the arbitrary point code as a parameter and induct on processed
prefixes with a bounded membership invariant. A universal coverage
assumption remains in the fixed context and is instantiated after the
finite program invariant has been proved; it is never an induction
formula.

\subsection{The final native transfer and integration boundary}

The actual upstream theorem is
\begin{lstlisting}
pra_iff_native :
  (Prf PRAAx Gamma phi) <->
  NPrf (Gamma.map NFm.ofPRA) (NFm.ofPRA phi)
\end{lstlisting}
in \code{ConciseIPRA.NativeBiInterpretation}~\cite{repo-native}.
It is applied after the object proof has been constructed. It does
not accept an ambient equality of denotations in place of its left-hand
proof object.

To integrate, replace the standalone kernel import with
\code{import ConciseIPRA.PR}, retain the new modules in their own
namespaces, complete the pending object targets, and use this bridge
for the final theorem. The standalone development was compiled.
A build against the full Mathlib-backed repository dependency closure
was not run, and no remote repository files were modified.
